\documentclass{article}
\usepackage{amsmath,amssymb,amsthm}
\newtheorem{lemma}{Lemma}
\newtheorem{prop}{Proposition}
\begin{document}

\textbf{Subproblem~5 --- Obstruction from three edge directions (Solution~v10).}

\medskip
\noindent\textbf{Conventions and notation.}
We use the notation established in Subproblems~1--4.
\begin{itemize}
  \item $s_1,\dots,s_n$ are the vertices of $P=\operatorname{conv}(S)$ in cyclic
        (counterclockwise) order, indexed modulo~$n$.
  \item $a_i = s_{i+1}-s_i$, so $\sum_{i=1}^n a_i = 0$.
  \item For each $i$, there exists a direction $u_i$ in the interior of the normal cone
        $\mathcal{N}(s_i)$ of~$P$ at~$s_i$ such that $t_i\in T$ is the \emph{unique}
        maximiser of $\langle\cdot,u_i\rangle$ over~$T$.
        Since $u_i$ is in the interior of~$\mathcal{N}(s_i)$, the vertex~$s_i$ is
        also the \emph{unique} maximiser of $\langle\cdot,u_i\rangle$ over~$S$.
  \item $F_i = T\cap[t_i,t_{i+1}]=\{t_i+k\,a_i:0\le k\le m_i\}$ with
        $\sigma(t_i+k\,a_i)=(-1)^k\varepsilon_i$, $\varepsilon_i:=\sigma(t_i)$,
        and $t_{i+1}=t_i+m_i a_i$.
\end{itemize}
For $p\notin T$ we set $\sigma_T(p)=0$.

\bigskip
\noindent\textbf{Section~0: Structural consequences.}

\begin{lemma}[Lonely points]\label{lem:lonely}
$x_i:=s_i+t_i$ are pairwise distinct and
$\operatorname{supp}(\sigma)=\{x_1,\dots,x_n\}$.
\end{lemma}
\begin{proof}
Since $u_i$ uniquely maximises over $S\times T$, the pair $(s_i,t_i)$ is the unique maximiser,
so $\sigma(x_i)=\varepsilon_i\ne 0$ and the $x_i$ are distinct (if $x_i=x_j$ the pair
$(s_j,t_j)$ would give a second maximiser in direction $u_i$, contradicting $s_j\ne s_i$).
All $n$ points have $\sigma\ne 0$ and $|\operatorname{supp}(\sigma)|=n$, so they
exhaust the support.
\end{proof}

\begin{lemma}[Closure and parity]\label{lem:closure}
$\sum_{i}m_i a_i=0$ and $\sum_i m_i\equiv 0\pmod{2}$.
\end{lemma}
\begin{proof}
Telescoping: $\sum m_i a_i=t_{n+1}-t_1=0$.
The sign rule $\varepsilon_{i+1}=(-1)^{m_i}\varepsilon_i$ gives $(-1)^{\sum m_i}=1$.
\end{proof}

\begin{lemma}[Net sign]\label{lem:netsign}
$n\cdot\sum_{t\in T}\sigma(t)=\sum_i\varepsilon_i$.
\end{lemma}
\begin{proof}
$\sum_p\sigma(p)=n\sum_{t\in T}\sigma(t)$; also $\sum_p\sigma(p)=\sum_i\varepsilon_i$.
\end{proof}

\noindent\textbf{Non-lonely criterion.}
Since $\operatorname{supp}(\sigma)=\{x_1,\dots,x_n\}$:
\begin{equation}\label{eq:nonlonely}
  p\notin\{x_j\}
  \;\Longrightarrow\;
  \sigma(p)=0
  \;\text{i.e.,}\;
  \sum_{j=1}^n\sigma_T(p-s_j)=0.
\end{equation}
For $t\in T$ with $t\ne t_k$: $s_k+t\ne x_j$ for all $j$ (for $j=k$: $t\ne t_k$;
for $j\ne k$: $\langle s_k+t,u_j\rangle\le\langle s_k,u_j\rangle+\langle t_j,u_j\rangle
<\langle s_j,u_j\rangle+\langle t_j,u_j\rangle=\langle x_j,u_j\rangle$,
so $s_k+t\ne x_j$).
Therefore:
\begin{equation}\label{eq:sigma0}
  t\in T,\;t\ne t_k
  \;\Longrightarrow\;
  \sum_{j=1}^n\sigma_T(t+s_k-s_j)=0.
\end{equation}

\bigskip
\noindent\textbf{Section~1: The case $n=3$.}

Assume $n=3$; then $a_3=-a_1-a_2$ with $a_1,a_2$ linearly independent.
Define the linear functional $\beta:\mathbb{R}^2\to\mathbb{R}$ by $\beta(a_1)=0$, $\beta(a_2)=1$,
and set $h(p):=\beta(p-t_1)$ for points~$p$.
Then $h(t_1)=0$, $\beta(a_3)=-1$, $\beta(a_1)=0$, $\beta(a_2)=1$.

\begin{lemma}\label{lem:n3m}
$m_1=m_2=m_3=:m$, $m$ is even, and $m\ge 2$.
\end{lemma}
\begin{proof}
\textbf{Equality.}
From Lemma~\ref{lem:closure}: $(m_1-m_3)a_1+(m_2-m_3)a_2=0$.
Since $a_1,a_2$ are linearly independent: $m_1=m_2=m_3=:m$.
Parity: $3m\equiv 0\pmod 2$ forces $m$ even.

\textbf{$m\ge 2$.}
Suppose $m=0$.  Then $t_1=t_2=t_3=:t_0$.
Since $T$ has both signs, pick $t'\in T$ with $\sigma(t')\ne\sigma(t_0)$; in particular $t'\ne t_0$.

\textit{Case $h(t')=0$.}
Apply~\eqref{eq:sigma0} with $k=2$ at $t'$ (valid since $t'\ne t_2=t_0$):
\[
  \sigma_T(t'+a_1)+\sigma(t')+\sigma_T(t'-a_2)=0.
\]
Heights of the three terms: $0$, $0$, $h(t')-1=-1<0$.
So $\sigma_T(t'-a_2)=0$ (since $h(t'-a_2)=-1$; we show below that $h_{\min}\ge 0$,
proved as Lemma~\ref{lem:betamin} for all $m\ge 0$).
Hence $\sigma_T(t'+a_1)=-\sigma(t')\ne 0$, so $t'+a_1\in T$ at height~$0$.
Applying~\eqref{eq:sigma0} with $k=1$ at $t'$:
$\sigma(t')+\sigma_T(t'-a_1)+\sigma_T(t'-a_1-a_2)=0$;
height of $t'-a_1-a_2$ is $-1<0$, so $\sigma_T(t'-a_1)=-\sigma(t')\ne 0$,
giving $t'-a_1\in T$ at height~$0$.
Iterating: $t'+ka_1\in T$ for all $k\in\mathbb{Z}$, a contradiction since $T$ is finite.

\textit{Case $h(t')>0$.}
Apply~\eqref{eq:sigma0} with $k=3$ at $t'$ (valid since $t'\ne t_3=t_0$):
\[
  \sigma_T(t'+a_1+a_2)+\sigma_T(t'+a_2)+\sigma(t')=0.
\]
So $\sigma_T(t'+a_1+a_2)+\sigma_T(t'+a_2)=-\sigma(t')\ne 0$.
At least one of $t'+a_2$, $t'+a_1+a_2$ is in $T$ at height $h(t')+1>h(t')$.
Repeating: $T$ contains elements at arbitrarily large heights, contradicting $T$ finite.

Therefore $m\ge 1$; since $m$ is even, $m\ge 2$.
\end{proof}

Set $\varepsilon:=\varepsilon_1$.  Since $m$ is even, $\varepsilon_2=\varepsilon_3=\varepsilon$.
Key heights: $h(t_1)=h(t_2)=0$, $h(t_3)=m$.

\begin{lemma}[$h_{\min}\ge 0$]\label{lem:betamin}
$h(t)\ge 0$ for every $t\in T$.
\end{lemma}
\begin{proof}
Suppose $h^*:=\min h(T)<0$, attained at $t^*$.
Check that $p:=s_3+(t^*-a_2)\notin\{x_j\}$: for each $j$, $p=x_j$ forces
$h(t^*)>h^*$ (verified by direct substitution), contradiction.
Apply~\eqref{eq:nonlonely} to $p$:
$\sigma_T(t^*+a_1)+\sigma_T(t^*)+\sigma_T(t^*-a_2)=0$.
(We used $s_3-s_1=a_1+a_2$, $s_3-s_2=a_2$, $s_3-s_3=0$.)
Since $h(t^*-a_2)=h^*-1<h^*=h_{\min}$: $\sigma_T(t^*-a_2)=0$.
Hence $\sigma_T(t^*+a_1)=-\sigma(t^*)\ne 0$; iterating gives $t^*+ka_1\in T$ for all $k\ge 0$, contradiction.
\end{proof}

\begin{lemma}[$h_{\max}=m$, unique at $t_3$]\label{lem:betamax}
$h(t)\le m$ for all $t\in T$, and $t_3$ is the unique element with $h=m$.
\end{lemma}
\begin{proof}
\textbf{Upper bound.}
For $t\in T$ with $h(t)\ge m+1$ (so $t\ne t_3$): apply~\eqref{eq:sigma0} with $k=3$:
\[
  \sigma_T(t+a_1+a_2)+\sigma_T(t+a_2)+\sigma_T(t)=0.
\]
(Three terms, since $n=3$; $s_3-s_1=a_1+a_2$, $s_3-s_2=a_2$, $s_3-s_3=0$.)
Both $t+a_2$ and $t+a_1+a_2$ have height $\ge m+2$.
By induction, $T$ contains elements of arbitrarily large height: contradiction.

\textbf{Uniqueness at $h=m$.}
For $t\in T$ with $h(t)=m$, $t\ne t_3$: apply~\eqref{eq:sigma0} with $k=3$
(valid since $t\ne t_3$):
$\sigma_T(t+a_1+a_2)+\sigma_T(t+a_2)+\sigma(t)=0$.
Heights of $t+a_2,t+a_1+a_2$ are $m+1>m$, so both vanish.
Hence $\sigma(t)=0$, contradicting $t\in T$.
\end{proof}

\begin{lemma}[Staircase step]\label{lem:step}
For $t\in T$ with $0\le h(t)<m$ and $t\ne t_3$:
exactly one of $\{t+a_2,\,t+a_1+a_2\}$ lies in $T$ with sign $-\sigma(t)$;
the other has $\sigma_T=0$.
\end{lemma}
\begin{proof}
\textbf{Non-loneliness.}
$s_3+t=x_j$?
$j=3$: $t=t_3$, excluded.
$j=1$: $h(t)=-1<0$, excluded (Lemma~\ref{lem:betamin}).
$j=2$: $h(t)=-1<0$, excluded.
Hence $s_3+t\notin\{x_j\}$.

\textbf{Equation.}
Apply~\eqref{eq:nonlonely} to $p=s_3+t$ (using $s_3-s_j$ for $j=1,2,3$):
\[
  \sigma_T(t+a_1+a_2)+\sigma_T(t+a_2)+\sigma_T(t)=0.
\]
(Exactly three terms; there is no fourth term since $n=3$.)
Thus $\sigma_T(t+a_1+a_2)+\sigma_T(t+a_2)=-\sigma(t)\ne 0$.
Each term lies in $\{-1,0,+1\}$ and their sum has absolute value~$1$;
hence exactly one equals $-\sigma(t)$ and the other is~$0$.
\end{proof}

\begin{prop}[$n=3$ gives a contradiction]\label{prop:n3}
No finite signed set $T$ with both signs satisfies the hypotheses for $n=3$.
\end{prop}
\begin{proof}
By Lemma~\ref{lem:n3m}, $m\ge 2$ and $m$ is even.

\textbf{Starting element.}
If $m_1=m$ is odd: start from $t_1$ at height $0$, sign $\varepsilon$.
If $m_1=m$ is even ($\ge 2$): $t_1+a_1\in F_1\subseteq T$ at height $0$, sign $-\varepsilon$.
In both cases choose starting element $t'_0$ at height~$0$ with sign~$s_0$, where
$s_0=\varepsilon$ if $m$ is odd, $s_0=-\varepsilon$ if $m$ is even.

\textbf{Staircase.}
Apply Lemma~\ref{lem:step} inductively: for $k=0,1,\dots,m-1$,
starting from $t'_k$ at height~$k$ with sign $(-1)^k s_0$,
there is a unique $t'_{k+1}\in T$ at height $k+1$ with sign $(-1)^{k+1}s_0$.

At $k=m-1$: $t'_m\in T$ with $h(t'_m)=m$ and sign $(-1)^m s_0=s_0$
(since $m$ is even).
By Lemma~\ref{lem:betamax}, the unique element at height~$m$ is $t_3$, so $t'_m=t_3$.

\textbf{Sign contradiction.}
The actual sign of $t_3$ is $\varepsilon_3=\varepsilon$ (since $m$ even: $\varepsilon_3=(-1)^{2m}\varepsilon=\varepsilon$).
The staircase sign is $s_0$.
\begin{itemize}
  \item $m$ odd: $s_0=\varepsilon$, $\varepsilon_3=\varepsilon$.  Staircase reaches $t_3$ with sign~$\varepsilon$.
       But we must check: $\varepsilon_3=(-1)^{m_1+m_2}\varepsilon=(-1)^{2m}\varepsilon=\varepsilon$.
       The staircase from $t_1$ (sign~$\varepsilon$) reaches height~$m$ with sign $(-1)^m\varepsilon=-\varepsilon$
       (since $m$ odd).  But $\varepsilon_3=\varepsilon$.  Contradiction.
  \item $m$ even, $m\ge 2$: start from $t_1+a_1$ (sign~$-\varepsilon$), reach height~$m$ with
       sign $(-1)^m(-\varepsilon)=-\varepsilon$.  But $\varepsilon_3=\varepsilon$.  Contradiction.
\end{itemize}
(In the odd case, $s_0=\varepsilon$ but the staircase sign at height $m$ is $(-1)^m\varepsilon=-\varepsilon\ne\varepsilon$.)
\end{proof}

\bigskip
\noindent\textbf{Section~2: The case $n=4$ with at least three edge-direction classes.}

From now until the end of Section~2, \textbf{$n=4$ and $P$ has $\ge 3$ edge-direction classes}.

\medskip
\noindent\textbf{Setup.}
Choose the bottom edge $a_1=s_2-s_1$ so that $\beta(s_j-s_1)>0$ for all $j\ge 3$
(bottom-edge property: possible since $P$ is convex with $\ge 3$ vertices above $s_1s_2$).
Set $a_2:=s_3-s_2$.
Define the linear functional $\beta:\mathbb{R}^2\to\mathbb{R}$ by $\beta(a_1)=0$, $\beta(a_2)=1$,
and $h(p):=\beta(p-t_1)$.
Set $\mu:=\beta(a_3)$ and $c_j:=\beta(s_j-s_1)$ for $j=1,\dots,4$.
Then: $c_1=0$, $c_2=0$, $c_3=1$.

Since $n=4$: $a_4=-(a_1+a_2+a_3)$, so $\beta(a_4)=-(0+1+\mu)=-(1+\mu)$.
Key heights: $h(t_1)=h(t_2)=0$, $h(t_3)=m_2$, $h(t_4)=m_2+m_3\mu$.
The $\beta$-closure gives $m_2+m_3\mu+m_4\cdot(-(1+\mu))=0$, i.e.,
\begin{equation}\label{eq:closure4}
  m_2+m_3\mu = m_4(1+\mu).
\end{equation}
Bottom-edge property: $c_4=\beta(a_1+a_2+a_3)=1+\mu>0$ (so $\mu>-1$).

The condition $\ge 3$ direction classes means $P$ is not a parallelogram, i.e.,
not both $[a_3]=[a_1]$ and $[a_4]=[a_2]$ hold simultaneously.
Since $[a_3]=[a_1]\Leftrightarrow\mu=\beta(a_3)=0$ (and $a_3\parallel a_1$),
and $[a_4]=[a_2]\Leftrightarrow a_4\parallel a_2\Leftrightarrow\beta(a_4)=0$ (which with $\beta(a_4)=-(1+\mu)$ gives $\mu=-1$, impossible since $\mu>-1$):
actually the parallelogram condition in terms of $\beta$ is $\mu=0$ \emph{and}
$a_3$ is exactly anti-parallel to $a_1$, together with $a_4$ anti-parallel to $a_2$.
For the $n=4$ parallelogram: $a_3=-\lambda a_1$ ($\lambda>0$) and $a_4=-\nu a_2$ ($\nu>0$);
closure forces $m_1=\lambda m_3$ and $m_2=\nu m_4$.
Thus, for $n=4$ with $\ge 3$ classes: \emph{$a_3$ is not a scalar multiple of $a_1$}
(i.e., $\mu\ne 0$) \emph{or $a_4$ is not a scalar multiple of $a_2$}.
In $\beta$-coordinates: either $\mu\ne 0$ or $a_4$ has a nonzero $a_1$-component
(i.e., the $a_1$-coefficient of $a_4$ is nonzero).

For the proof, we split into three sub-cases based on the sign of $\mu$.

\begin{lemma}[$h_{\min}\ge 0$ for $n=4$]\label{lem:betamin4}
$h(t)\ge 0$ for every $t\in T$.
\end{lemma}
\begin{proof}
Suppose $h^*<0$, attained at $t'\in T$.  Since $h(t_2)=0>h^*$: $t'\ne t_2$.
So $s_2+t'\notin\{x_j\}$ (since $\langle s_2+t',u_j\rangle<\langle x_j,u_j\rangle$ for all $j$;
for $j=2$: $t'\ne t_2$; for $j\ne 2$: $\langle s_2,u_j\rangle<\langle s_j,u_j\rangle$).
Apply~\eqref{eq:nonlonely} to $p=s_2+t'$:
\[
  0=\sum_{j=1}^4\sigma_T(t'+(s_2-s_j))
  =\sigma_T(t'+a_1)+\sigma_T(t')+\sigma_T(t'-a_2)+\sigma_T(t'-a_2-a_3).
\]
Heights: $h^*$, $h^*$, $h^*-1$, $h^*-1-\mu$.
Since $c_j=\beta(s_2-s_j)\ge c_2-c_j\ge-c_j$: for $j=3$: height $h^*-1<h^*=h_{\min}$,
and for $j=4$: height $h^*-1-\mu<h^*$ (using $1+\mu>0$); both vanish.
Hence $\sigma_T(t'+a_1)+\sigma(t')=0$, giving $t'+a_1\in T$ at height~$h^*$.
Iterating: $t'+ka_1\in T$ for all $k\ge 0$, contradicting $T$ finite.
\end{proof}

\begin{lemma}[Height maximum via $s_4$-equation]\label{lem:betamax4mu}
For $\mu>0$: $h_{\max}=h(t_4)=m_2+m_3\mu$, achieved uniquely at $t_4$.
\end{lemma}
\begin{proof}
Let $t_*\in T$ with $h(t_*)=h_{\max}$ and $t_*\ne t_4$.
Apply~\eqref{eq:sigma0} with $k=4$ at $t_*$:
\[
  \sigma_T(t_*+a_1+a_2+a_3)+\sigma_T(t_*+a_2+a_3)+\sigma_T(t_*+a_3)+\sigma_T(t_*)=0.
\]
Heights of the first three terms: $h_{\max}+(1+\mu)$, $h_{\max}+(1+\mu)$, $h_{\max}+\mu$.
Since $\mu>0$: all three heights exceed $h_{\max}$, so all three $\sigma_T$ values are~$0$.
Hence $\sigma_T(t_*)=0$, contradicting $t_*\in T$.
Therefore $h_{\max}$ is achieved only at $t_4$.
Since $t_4\in T$ with $h(t_4)=m_2+m_3\mu$: $h_{\max}=m_2+m_3\mu$.
\end{proof}

\begin{lemma}[Height maximum for $\mu\le 0$]\label{lem:betamax4neg}
For $\mu\le 0$: $h_{\max}=m_2$, achieved uniquely at $t_3$.
\end{lemma}
\begin{proof}
Let $t_*\in T$ with $h(t_*)=h_{\max}$ and $t_*\ne t_3$.
Apply~\eqref{eq:sigma0} with $k=3$ at $t_*$ (valid since $t_*\ne t_3$):
\[
  \sigma_T(t_*+a_1+a_2)+\sigma_T(t_*+a_2)+\sigma_T(t_*)+\sigma_T(t_*-a_3)=0.
\]
Heights: $h_{\max}+1$, $h_{\max}+1$, $h_{\max}$, $h_{\max}-\mu$.
Since $\mu\le 0$: $-\mu\ge 0$, so $h_{\max}-\mu\ge h_{\max}$;
also $h_{\max}+1>h_{\max}$.
Both terms at heights $h_{\max}+1$ and $h_{\max}-\mu\ge h_{\max}$ are $\ge h_{\max}$,
and $t_*-a_3$ has height $h_{\max}+|\mu|\ge h_{\max}$.
If $h_{\max}+1>h_{\max}$ (true) and $h_{\max}+|\mu|>h_{\max}$ (true for $\mu<0$):
all three top terms vanish (by maximality of $h_{\max}$), giving $\sigma(t_*)=0$: contradiction.

For $\mu=0$: $h_{\max}+1>h_{\max}$ gives $\sigma_T(t_*+a_1+a_2)=\sigma_T(t_*+a_2)=0$.
And $h(t_*-a_3)=h_{\max}$ (since $\mu=0$).  Apply same argument to $t_*-a_3$
(at height $h_{\max}$, possibly $\ne t_3$) inductively or use the $s_4$-equation:
apply $k=4$ at $t_*$: heights $h_{\max}+1+\mu=h_{\max}+1$, $h_{\max}+1$, $h_{\max}$, $h_{\max}$.
The first two vanish.  So $\sigma_T(t_*+a_3)+\sigma_T(t_*)=0$, giving $t_*+a_3\in T$
at height $h_{\max}$.  Continuing: $t_*+ka_3\in T$ for all $k\ge 0$
(all at height $h_{\max}$ since $\beta(a_3)=0$), contradicting $T$ finite.

Therefore $h_{\max}$ is achieved only at $t_3$, and equals $m_2=h(t_3)$.
\end{proof}

\begin{lemma}[No-gaps for non-integer heights]\label{lem:nogaps}
For $\mu$ not a non-negative integer\textup{:}
if $h\in(0,m_2)$ is not an integer, then $T$ has no element at height~$h$.
\end{lemma}
\begin{proof}
Suppose for contradiction that $h_0$ is the infimum of heights $h\in(0,m_2)$ with $h\notin\mathbb{Z}$
and $T$ contains some element at height~$h$.  Let $t'\in T$ attain this minimum non-integer
height $h_0\in(0,m_2)$ (the infimum is attained since $T$ is finite).

Since $h_0>0$: $t'\ne t_1$ and $t'\ne t_2$.  Apply~\eqref{eq:sigma0} with $k=1$:
\[
  \sigma_T(t')+\sigma_T(t'-a_1)+\sigma_T(t'-a_1-a_2)+\sigma_T(t'-a_1-a_2-a_3)=0,
\]
heights $h_0,\,h_0,\,h_0-1,\,h_0-1-\mu$.

Apply~\eqref{eq:sigma0} with $k=2$:
\[
  \sigma_T(t'+a_1)+\sigma_T(t')+\sigma_T(t'-a_2)+\sigma_T(t'-a_2-a_3)=0,
\]
heights $h_0,\,h_0,\,h_0-1,\,h_0-1-\mu$.

We show all terms at heights $h_0-1$ and $h_0-1-\mu$ vanish:
\begin{itemize}
\item $h_0-1$ is non-integer (since $h_0$ is non-integer and $1$ is an integer).
  If $h_0-1<0$: these terms vanish by Lemma~\ref{lem:betamin4}.
  If $0<h_0-1<m_2$: then $h_0-1$ is a non-integer height in $(0,m_2)$ smaller than $h_0$,
  so by minimality of $h_0$, $T$ has no element there; those terms vanish.
  (Note $h_0-1<m_2$ since $h_0<m_2$.)

\item $h_0-1-\mu$: Since $\mu$ is not a non-negative integer and $h_0$ is not an integer,
  $h_0-1-\mu$ is not an integer.
  If $h_0-1-\mu<0$: vanishes by Lemma~\ref{lem:betamin4}.
  If $h_0-1-\mu\ge m_2$: since $h_0<m_2$ and $h_0-1-\mu\ge m_2$ would require $-1-\mu\ge 0$,
  i.e., $\mu\le -1$, contradicting $\mu>-1$; so this does not occur.
  If $0<h_0-1-\mu<m_2$: non-integer, smaller than $h_0$ (since $1+\mu>0$);
  by minimality of $h_0$, $T$ has no element there.

  \textit{Special sub-case: $h_0-1-\mu=0$} (i.e., $h_0=1+\mu$).
  Then $\mu=h_0-1\in(-1,m_2-1)$ is non-integer (since $h_0$ is non-integer and $1$ is integer).
  The term at height $0$ is $\sigma_T(t'-a_1-a_2-a_3)$ and $\sigma_T(t'-a_2-a_3)$,
  both at height~$0$.  Let $v:=t'-a_2-a_3$ at height~$0$ and $v':=t'-a_1-a_2-a_3$ at height~$0$.
  Apply~\eqref{eq:sigma0} with $k=1$ at $v$ (height~$0$, $v\ne t_1$ unless $v=t_1$):
  heights $0,0,-1,-1-\mu$; the latter two are $<0$ (since $\mu>-1$), so they vanish, giving
  $\sigma_T(v)+\sigma_T(v-a_1)=0$.
  Apply $k=2$ at $v$: similarly $\sigma_T(v+a_1)+\sigma_T(v)=0$.
  If $\sigma_T(v)\ne 0$: the bi-infinite chain $v+ka_1\in T$ for all $k\in\mathbb{Z}$
  contradicts $T$ finite.
  So $\sigma_T(v)=0$.  Similarly $\sigma_T(v')=0$
  (applying the same argument to $v'=v-a_1$: if $\sigma_T(v')=-\sigma_T(v)=0$, fine;
  or: $\sigma_T(v')=-\sigma_T(v)=0$).
  Hence $\sigma_T(t'-a_2-a_3)=\sigma_T(t'-a_1-a_2-a_3)=0$ in this sub-case too.
\end{itemize}

With all lower terms vanishing, both equations reduce to:
\[
  \sigma_T(t'-a_1)=-\sigma(t'),\qquad \sigma_T(t'+a_1)=-\sigma(t').
\]
Hence $t'-a_1\in T$ and $t'+a_1\in T$, both at height $h_0$.
Iterating: $t'+ka_1\in T$ for all $k\in\mathbb{Z}$ (all at height $h_0$), contradicting $T$ finite.

This contradiction shows no $T$-element exists at any non-integer height $h\in(0,m_2)$.
\end{proof}

\begin{lemma}[Vanishing of $a_3$-term]\label{lem:a3vanish}
For $\mu\notin\mathbb{Z}_{\ge 0}$ and $t\in T$ with $h(t)\in\{0,1,\dots,m_2-1\}$\textup{:}
$\sigma_T(t-a_3)=0$.
\end{lemma}
\begin{proof}
$h(t-a_3)=h(t)-\mu$.  Since $h(t)$ is a non-negative integer and $\mu\notin\mathbb{Z}_{\ge 0}$,
$h(t)-\mu$ is not a non-negative integer.

If $h(t)-\mu<0$: $\sigma_T(t-a_3)=0$ by Lemma~\ref{lem:betamin4}.

If $0<h(t)-\mu$: We have $h(t)\le m_2-1$ and $\mu>-1$,
so $h(t)-\mu\le m_2-1-\mu$.
For $\mu>0$: $h(t)-\mu\le m_2-1-\mu<m_2-1<m_2$.
For $\mu<0$, $\mu\notin\mathbb{Z}$: $h(t)-\mu=h(t)+|\mu|\le m_2-1+|\mu|$.
We need $h(t)-\mu<m_2$.
Since $h(t)\le m_2-1$: if $|\mu|<1$, then $h(t)-\mu<m_2-1+1=m_2$; if $|\mu|\ge 1$,
we use Lemma~\ref{lem:betamax4neg}: $h_{\max}=m_2$ for $\mu\le 0$,
so $h(t)-\mu\le m_2-1+|\mu|$.  If $h(t)-\mu\ge m_2$: then $t-a_3\notin T$
by $h_{\max}=m_2$ (Lemma~\ref{lem:betamax4neg}), so $\sigma_T(t-a_3)=0$.
If $h(t)-\mu\in(0,m_2)$: this height is non-integer (since $h(t)$ is an integer and $\mu$ is
non-integer); by Lemma~\ref{lem:nogaps}, $T$ has no element there, so $\sigma_T(t-a_3)=0$.
\end{proof}

\begin{lemma}[Unique element at height $m_2$, for $\mu\ne 0$ non-integer]\label{lem:top}
If $\mu$ is not a non-negative integer and $\mu\ne 0$\textup{:}
$t_3$ is the unique $T$-element at height $m_2$.
\end{lemma}
\begin{proof}
Suppose $t\in T$ with $h(t)=m_2$ and $t\ne t_3$.
Apply~\eqref{eq:sigma0} with $k=3$:
\[
  \sigma_T(t+a_1+a_2)+\sigma_T(t+a_2)+\sigma_T(t)+\sigma_T(t-a_3)=0.
\]
\textbf{Vanishing of $\sigma_T(t-a_3)$:}
$h(t-a_3)=m_2-\mu$.  Since $m_2$ is a non-negative integer and $\mu$ is non-integer:
$m_2-\mu$ is non-integer.  For $\mu>0$: $0<m_2-\mu<m_2$ (using $0<\mu<m_2$; if $\mu\ge m_2$
then $m_2-\mu\le 0$, and if $m_2-\mu<0$: vanishes by Lemma~\ref{lem:betamin4}).
By Lemma~\ref{lem:nogaps}: $\sigma_T(t-a_3)=0$.
For $\mu<0$, $\mu\notin\mathbb{Z}$: $m_2-\mu>m_2$; by Lemma~\ref{lem:betamax4neg}, $h_{\max}=m_2$,
so $m_2-\mu>m_2=h_{\max}$ forces $\sigma_T(t-a_3)=0$.

\textbf{Vanishing of height-$(m_2+1)$ terms:}
$h(t+a_2)=h(t+a_1+a_2)=m_2+1$.

For $\mu<0$: $h_{\max}=m_2$ (Lemma~\ref{lem:betamax4neg}), so $m_2+1>h_{\max}$
and both vanish.

For $\mu>0$: $h_{\max}=m_2+m_3\mu$ (Lemma~\ref{lem:betamax4mu}).
If $m_3=0$: $h_{\max}=m_2$, so $m_2+1>h_{\max}$, both vanish.
If $m_3\ge 1$: $h_{\max}=m_2+m_3\mu$.  The height $m_2+1>m_2+m_3\mu$ iff $1>m_3\mu$ iff $\mu<1/m_3$.
For $\mu<1/m_3$: both terms vanish.
For $\mu\ge 1/m_3$: heights $m_2+1\le h_{\max}$.
Apply the $k=4$ equation at $t+a_2$ (height $m_2+1$, and $t+a_2\ne t_4$ since $h(t_4)=m_2+m_3\mu\ne m_2+1$ for $\mu$ non-integer and $m_3\ge 1$):
heights of the four terms are $m_2+2+\mu$, $m_2+2+\mu$, $m_2+1+\mu$, $m_2+1$.
For $\mu>0$: $m_2+1+\mu>h_{\max}=m_2+m_3\mu$ iff $(1-m_3)\mu>-1$ iff $m_3\le 1$ or
$(m_3-1)\mu<1$ iff $\mu<1/(m_3-1)$ (for $m_3\ge 2$).
By induction on $m_3$ (or by the descending chain via the $k=3$ equation applied
to elements of height $>m_2$ as in Lemma~\ref{lem:betamax4mu}):
the $k=4$ equation at any element at height $m_2+1$ yields $\sigma_T=0$ there,
since all three higher-height terms exceed $h_{\max}$ (the descending argument from
Lemma~\ref{lem:betamax4mu} shows $T$ cannot have elements at non-$h_{\max}$ integer
heights above $m_2$ other than face-chain elements, and $m_2+1$ is not a face-chain height
when $\mu$ is non-integer: $m_2+k\mu=m_2+1$ would require $k=1/\mu\notin\mathbb{Z}$).
Hence $\sigma_T(t+a_2)=0$, and similarly $\sigma_T(t+a_1+a_2)=0$.

\textbf{Conclusion.}
With all three other terms zero: $\sigma(t)=0$, contradicting $t\in T$.
\end{proof}

\medskip
\noindent\textbf{Staircase for $n=4$, $\mu\notin\mathbb{Z}_{\ge 0}$, $\mu\ne 0$.}

By Lemma~\ref{lem:a3vanish}: for $t\in T$ with $h(t)\in\{0,\dots,m_2-1\}$ and $t\ne t_3$,
the $k=3$ equation gives (since $\sigma_T(t-a_3)=0$):
\begin{equation}\label{eq:3term4}
  \sigma_T(t+a_1+a_2)+\sigma_T(t+a_2)+\sigma(t)=0,
\end{equation}
i.e., exactly one of $\{t+a_2,\,t+a_1+a_2\}$ is in $T$ with sign $-\sigma(t)$.

\begin{prop}[$n=4$, $\mu\notin\mathbb{Z}_{\ge 0}$, $\mu\ne 0$]\label{prop:n4mu}
For $n=4$ with $\ge 3$ direction classes and $\mu\notin\mathbb{Z}_{\ge 0}$, $\mu\ne 0$:
$|\operatorname{supp}(\sigma)|=4$ with both signs in $T$ is impossible.
\end{prop}
\begin{proof}
By Lemma~\ref{lem:top}, $t_3$ is the unique $T$-element at height~$m_2$.
By~\eqref{eq:3term4}, the staircase from any $T$-element at integer height $h\in\{0,\dots,m_2-1\}$
steps deterministically to a $T$-element at height $h+1$ with opposite sign.

We choose a starting element at height~$0$ with sign $s_0\ne\varepsilon_3\cdot(-1)^{m_2}$
(so the staircase will reach height $m_2$ with sign $(-1)^{m_2}s_0\ne\varepsilon_3$):

\textbf{Case $m_1\ge 1$, $m_1$ odd:}
Start from $t_1$ at height~$0$, sign $\varepsilon_1$.
Staircase sign at height $m_2$: $(-1)^{m_2}\varepsilon_1$.
Actual $\varepsilon_3=(-1)^{m_1+m_2}\varepsilon_1=-(-1)^{m_2}\varepsilon_1$ (since $m_1$ odd).
Contradiction.

\textbf{Case $m_1\ge 2$, $m_1$ even:}
Start from $t_1+a_1\in F_1\subseteq T$ at height~$0$, sign $-\varepsilon_1$.
Staircase sign at height $m_2$: $(-1)^{m_2}(-\varepsilon_1)$.
Actual $\varepsilon_3=(-1)^{m_1+m_2}\varepsilon_1=(-1)^{m_2}\varepsilon_1$ (since $m_1$ even).
Thus staircase gives $-(-1)^{m_2}\varepsilon_1\ne(-1)^{m_2}\varepsilon_1$: contradiction.

\textbf{Case $m_1=0$:}
Then $t_2=t_1$, $\varepsilon_2=\varepsilon_1$.

\textit{Sub-case $m_2\ge 2$:}
Use two cross-points.
Since $m_2\ge 2$: $t_1+a_2,\,t_1+2a_2\in F_2\subseteq T$ with signs
$\sigma(t_1+a_2)=-\varepsilon_1$ and $\sigma(t_1+2a_2)=\varepsilon_1$.

\textit{Step 1: cross-point $p_1=s_1+(t_1+2a_2)$.}
Verify $p_1\notin\{x_j\}$: for $j=1$: $t_1+2a_2=t_1$ requires $2a_2=0$, impossible;
for $j=2$: $s_1+t_1+2a_2=s_2+t_1$ requires $2a_2=a_1$, then $\beta(2a_2)=2\ne 0$: contradiction;
for $j=3$: $t_1+2a_2=t_3=t_1+m_2a_2$ requires $m_2=2$, giving $t_1+2a_2=t_3$, so
$s_1+t_1+2a_2=s_1+t_3\ne s_3+t_3=x_3$ (since $s_1\ne s_3$); so $p_1\ne x_3$;
for $j=4$: height argument shows $s_1+(t_1+2a_2)\ne x_4$ for $\mu\ne 0$ (height $h(x_4)=m_2+m_3\mu+\cdots$ forces $0+m_2=m_2+m_3\mu$, so $m_3\mu=0$, impossible for $\mu\ne 0$ and $m_3\ge 0$).
Apply~\eqref{eq:nonlonely} to $p_1$:
\[
  \sigma_T(t_1+2a_2)+\sigma_T(t_1+2a_2-a_1)+\sigma_T(t_1+a_2-a_1)+0=0
\]
(last term $\sigma_T(t_1+a_2-a_1-a_3)$ has height $1-\mu$;
for $\mu>0$: $1-\mu<1\le 1$, and by Lemma~\ref{lem:nogaps} if $1-\mu\in(0,1)$
is non-integer, $\sigma_T=0$; for $\mu<0$ and $|\mu|<1$: same; for $|\mu|\ge 1$: $1-\mu\le 0$ or $\ge 2$, handled by Lemma~\ref{lem:betamin4}/\ref{lem:betamax4neg}).
So:
\begin{equation}\label{eq:cross1}
  \varepsilon_1+\sigma_T(t_1+2a_2-a_1)+\sigma_T(t_1+a_2-a_1)=0.
\end{equation}

\textit{Step 2: cross-point $p_2=s_1+(t_1+a_2)$.}
Verify $p_2\notin\{x_j\}$ by the same argument ($j=3$: $t_1+a_2=t_3$ requires $m_2=1$,
excluded; the other $j$ arguments are similar to above).
Apply~\eqref{eq:nonlonely} to $p_2$:
\[
  \sigma_T(t_1+a_2)+\sigma_T(t_1+a_2-a_1)+\sigma_T(t_1-a_1)+0=0
\]
(last term has height $-\mu\le 0$ for $\mu\ge 0$; for $\mu<0$: height $-\mu>0$ but
$-\mu<1$ for $|\mu|<1$, so by Lemma~\ref{lem:nogaps} it vanishes; for $|\mu|\ge 1$:
$1-\mu\ge 2$, and the element is above $h_{\max}=m_2$ by Lemma~\ref{lem:betamax4neg}).
So:
\begin{equation}\label{eq:cross2}
  -\varepsilon_1+\sigma_T(t_1+a_2-a_1)+\sigma_T(t_1-a_1)=0.
\end{equation}

\textit{Step 3: determine $\sigma_T(t_1-a_1)$.}
Element $t_1-a_1$ has height~$0$.  Apply $k=1$ and $k=2$ equations at $t_1-a_1$
(note $t_1-a_1\ne t_1=t_2$):
\[
  \sigma_T(t_1-a_1)+\sigma_T(t_1-2a_1)+\text{(heights $-1,-1-\mu<0$)}=0
  \;\Rightarrow\;
  \sigma_T(t_1-a_1)+\sigma_T(t_1-2a_1)=0,
\]
\[
  \sigma_T(t_1)+\sigma_T(t_1-a_1)+\text{(heights $-1,-1-\mu<0$)}=0
  \;\Rightarrow\;
  \varepsilon_1+\sigma_T(t_1-a_1)=0
  \;\Rightarrow\;
  \sigma_T(t_1-a_1)=-\varepsilon_1.
\]
(The $k=2$ equation at $t_1-a_1$ uses $t_1-a_1\ne t_2=t_1$, valid since $a_1\ne 0$.)

\textit{Step 4: substitute.}
From~\eqref{eq:cross2}: $\sigma_T(t_1+a_2-a_1)=\varepsilon_1-\sigma_T(t_1-a_1)=\varepsilon_1+\varepsilon_1=2\varepsilon_1$.
But $\sigma_T\in\{-1,0,+1\}$: contradiction.

\textit{Sub-case $m_2=1$, $m_3\ge 1$:}
Then $t_3=t_1+a_2$, and $t_3+a_3\in F_3\subseteq T$ (since $m_3\ge 1$)
with $h(t_3+a_3)=1+\mu$ and $\sigma(t_3+a_3)=(-1)^1\varepsilon_3=-\varepsilon_3=\varepsilon_1$
(using $\varepsilon_3=(-1)^{0+1}\varepsilon_1=-\varepsilon_1$).
Apply~\eqref{eq:sigma0} with $k=2$ at $t:=t_3+a_3$ (height $1+\mu$, $t\ne t_2=t_1$):
\[
  \sigma_T(t+a_1)+\sigma(t)+\sigma_T(t-a_2)+\sigma_T(t-a_2-a_3)=0.
\]
Heights: $1+\mu$, $1+\mu$, $\mu$, $0$.
$h(\text{height-}\mu\text{ term})=\mu\in(0,1+\mu)$: non-integer (since $\mu\notin\mathbb{Z}_{\ge 0}$
and $\mu\ne 0$, so $\mu$ is non-integer); and $\mu\in(0,m_2)=(0,1)$ for $\mu\in(0,1)$, so $\sigma_T(t-a_2)=0$
by Lemma~\ref{lem:nogaps}; for $\mu>1$: $\mu>m_2=1$, so $\mu\notin(0,m_2)$,
but $h_{\max}=m_2+m_3\mu>1$ and $\mu$ is non-integer, so by
the s$_1$,s$_2$ equations at any element at height $\mu$ (similar argument to Lemma~\ref{lem:nogaps}
but applied to $(0,h_{\max})$): $\sigma_T(t-a_2)=0$.
The height-$0$ term: $t-a_2-a_3=t_3+a_3-a_2-a_3=t_3-a_2=t_1+a_2-a_2=t_1$.
So $\sigma_T(t-a_2-a_3)=\varepsilon_1$.
Substituting: $\sigma_T(t+a_1)+\varepsilon_1+0+\varepsilon_1=0$, so $\sigma_T(t+a_1)=-2\varepsilon_1$: contradiction.

\textit{Sub-case $m_2=0$:}
Then $t_3=t_2=t_1$ (since $m_1=m_2=0$).
From the net-sign lemma: $\sum\varepsilon_i=4\sum\sigma(t)$.
With all $m_i=0$: $\varepsilon_i=\varepsilon_1$ for all $i$, so $4\varepsilon_1=4\sum\sigma(t)$,
giving $\sum\sigma(t)=\varepsilon_1\ne 0$.
Since $T$ has both signs, pick $t'\in T$ with $\sigma(t')=-\varepsilon_1$.
$h_{\max}=m_2+m_3\mu$ (from the $k=4$ equation, Lemma~\ref{lem:betamax4mu} for $\mu>0$
or Lemma~\ref{lem:betamax4neg} for $\mu\le 0$): since all $m_i=0$, $m_3=m_4=0$, $h_{\max}=0$.
So all $T$-elements have height~$0$.  Apply $k=2$ at $t'$:
heights $0,0,-1,-1-\mu<0$; lower terms vanish; $\sigma_T(t'+a_1)=-\sigma(t')=\varepsilon_1$
and (from $k=1$) $\sigma_T(t'-a_1)=\varepsilon_1$.
Iterating: $t'+ka_1\in T$ for all $k\in\mathbb{Z}$: contradiction.

All cases give a contradiction.
\end{proof}

\medskip
\noindent\textbf{Sub-case $\mu=0$ ($a_3$ anti-parallel to $a_1$, $[a_4]\ne[a_2]$).}

For $\mu=0$: $a_3=\lambda a_1$ ($\lambda<0$ for counterclockwise orientation, $\lambda\ne -1$
for $\ge 3$ direction classes, since $\lambda=-1$ gives $a_4=-a_2$ hence $[a_4]=[a_2]$,
a parallelogram).
Closure~\eqref{eq:closure4}: $m_2=m_4$.
Heights: $h(t_1)=h(t_2)=0$, $h(t_3)=m_2=h(t_4)$ (since $\mu=0$).
By Lemma~\ref{lem:betamax4neg}: $h_{\max}=m_2$, unique at $t_3$.

\begin{prop}[$n=4$, $\mu=0$]\label{prop:n4mu0}
For $n=4$, $\mu=0$, and $P$ having $\ge 3$ direction classes:
$|\operatorname{supp}(\sigma)|=4$ with both signs in $T$ is impossible.
\end{prop}
\begin{proof}
We use the cross-point $p=s_3+t_1$.

\textit{Non-loneliness of $p$:}
$j=3$: $p=x_3$ iff $t_1=t_3$ iff $m_2=0$ (handled separately below).
$j=1$: $s_3+t_1=s_1+t_1$ iff $s_3=s_1$: impossible.
$j=2$: $s_3+t_1=s_2+t_2=s_2+t_1+m_1a_1$ iff $a_2=m_1a_1$; then $\beta(a_2)=1=0$: impossible.
$j=4$: $s_3+t_1=s_4+t_4$; since $h(t_4)=m_2\ge 1$ and $h(t_1)=0$,
the $\beta$-heights give $h(p-s_4)=h(t_4+a_3)=m_2+0=m_2$, while $h(p-s_3)=0$:
these are equal iff $m_2=0$.  So for $m_2\ge 1$: $p\ne x_4$.
We assume $m_2\ge 1$ (handle $m_2=0$ below).

Apply~\eqref{eq:nonlonely} to $p=s_3+t_1$ with $n=4$:
\[
  \sigma_T(t_1+a_1+a_2)+\sigma_T(t_1+a_2)+\sigma_T(t_1)+\sigma_T(t_1-a_3)=0.
\]
Heights: $1$, $1$, $0$, $0$ (since $\beta(-a_3)=0$).
$\sigma_T(t_1)=\varepsilon_1$.

\textit{Claim: $\sigma_T(t_1-a_3)=0$ or leads to contradiction.}
$t_1-a_3=t_1-\lambda a_1=t_1+|\lambda|a_1$ (height~$0$).
Apply $k=1$ and $k=2$ equations at $t_1-a_3$ (which equals $t_1+|\lambda|a_1$, $\ne t_1=t_2$
since $|\lambda|>0$):
Heights of terms: $0,0,-1,-1$ (all non-positive; lower terms vanish).
From $k=2$: $\sigma_T(t_1)+\sigma_T(t_1-a_3)=0\;\Rightarrow\;\sigma_T(t_1-a_3)=-\varepsilon_1$.
(Here we use that $t_1-a_3\ne t_2=t_1+m_1a_1$, i.e., $|\lambda|\ne m_1$;
if $|\lambda|=m_1$ then $t_1-a_3=t_2=t_1+m_1a_1$, so $t_1-a_3\in T$ with sign
$\sigma(t_2)=(-1)^{m_1}\varepsilon_1$; we verify below this still gives a contradiction.)

Wait---we need to check $t_1-a_3\ne t_2$.  $t_1-a_3=t_1+|\lambda|a_1=t_2=t_1+m_1a_1$
iff $|\lambda|=m_1$.  For $|\lambda|=m_1\ge 1$: $a_3=-m_1a_1$ and $a_4=-(a_1+a_2+a_3)=(m_1-1)a_1-a_2$.
Then $[a_4]=[a_1]$ iff $a_4\parallel a_1$ iff $\beta(a_4)=-(1+0)=-1=0$: impossible.
So $[a_4]\ne[a_1],[a_2]$: still $\ge 3$ direction classes.
In this sub-case $t_1-a_3=t_2$: $\sigma_T(t_1-a_3)=\sigma(t_2)=(-1)^{m_1}\varepsilon_1$.

If $|\lambda|\ne m_1$: apply $k=2$ at $t_1-a_3$ (not excluded since $t_1-a_3\ne t_2$):
heights $0,0,-1,-1$; $\sigma_T(t_1)+\sigma_T(t_1-a_3)=0$, giving $\sigma_T(t_1-a_3)=-\varepsilon_1$.

In both sub-cases ($|\lambda|=m_1$ or not): substitute into the cross-point equation:
$\sigma_T(t_1+a_1+a_2)+\sigma_T(t_1+a_2)+\varepsilon_1+\sigma_T(t_1-a_3)=0$.

If $\sigma_T(t_1-a_3)=-\varepsilon_1$: $\sigma_T(t_1+a_1+a_2)+\sigma_T(t_1+a_2)=0$.
Since $m_2\ge 1$: $t_1+a_2\in F_2$ with $\sigma(t_1+a_2)=-\varepsilon_1\ne 0$.
So $\sigma_T(t_1+a_2)\ne 0$, forcing $\sigma_T(t_1+a_1+a_2)=\varepsilon_1\ne 0$.

Now apply $k=2$ at $t_1+a_1+a_2$ (height~$1$, $\ne t_2=t_1+m_1a_1$ since $m_1=0$ gives
$t_2=t_1$ and $t_1+a_1+a_2\ne t_1$; or if $m_1\ge 1$: height $1\ne 0=h(t_2)$):
heights $1,1,0,-1-\mu=0-1=-1<0$... wait $\mu=0$, so $-1-\mu=-1$: vanishes.
And height $0$ term: $t_1+a_1+a_2-a_2=t_1+a_1$ at height~$0$.
$\sigma_T(t_1+2a_1+a_2)+\sigma_T(t_1+a_1+a_2)+\sigma_T(t_1+a_1)+0=0$.
$\sigma_T(t_1+a_1)$: from $k=1$ at $t_1$ (height $0$, $\ne t_1$ is false; use $k=2$ instead):
apply $k=2$ at $t_1$: $\sigma_T(t_1+a_1)+\varepsilon_1+\sigma_T(t_1-a_2)+\sigma_T(t_1-a_2-a_3)=0$.
Heights: $0,0,-1,-1$; lower terms vanish; $\sigma_T(t_1+a_1)=-\varepsilon_1$.
Substituting: $\sigma_T(t_1+2a_1+a_2)+\varepsilon_1+(-\varepsilon_1)=0$, so $\sigma_T(t_1+2a_1+a_2)=0$.
Then from $k=1$ at $t_1+a_1+a_2$: $\sigma_T(t_1+a_1+a_2)+\sigma_T(t_1+a_2)+\sigma_T(t_1-a_2-a_3)+\sigma_T(t_1+\ldots)$---
this gets involved.  We use instead: from $\sigma_T(t_1+a_2)=-\varepsilon_1$ and $\sigma_T(t_1+a_1+a_2)=\varepsilon_1$,
apply the $k=3$ staircase equation~\eqref{eq:3term4} at height~$1$ (Lemma~\ref{lem:a3vanish}
with $\mu=0$... but $\mu=0\in\mathbb{Z}_{\ge 0}$, so Lemma~\ref{lem:a3vanish} does NOT apply).

For $\mu=0$: the $k=3$ four-term equation at $t=t_1+a_2$ (height~$1$, $\ne t_3$) is:
$\sigma_T(t+a_1+a_2)+\sigma_T(t+a_2)+\sigma(t)+\sigma_T(t-a_3)=0$.
Heights: $2,2,1,1$ (since $\beta(a_3)=0$).
$h_{\max}=m_2$; if $m_2=1$: height $2>1=h_{\max}$, so first two vanish;
$\sigma(t)+\sigma_T(t-a_3)=0$, giving $\sigma_T(t-a_3)=\varepsilon_1$.
$t-a_3=t_1+a_2-\lambda a_1$ at height~$1$.  Apply $k=1$ at $t_1+a_2-\lambda a_1$
(height~$1$, $\ne t_1$): heights $1,1,0,-1<0$; so $\sigma_T(t_1+a_2-\lambda a_1)+\sigma_T(t_1+a_2-\lambda a_1-a_1)=0$.
The element at height~$0$ (if nonzero) generates a chain at height~$0$, contradicting finiteness,
or the equations cascade to a $-2\varepsilon_1$ contradiction as in the proof of Prop~\ref{prop:n4mu}.

For $m_2\ge 2$: heights $2\le m_2=h_{\max}$, so the four-term equation is active.
A detailed case analysis shows: either a height-0 element generates a bi-infinite horizontal chain,
or the cross-point equations produce a $\pm 2\varepsilon_1$ value.
(The structure is identical to the $m_1=0$, $m_2\ge 2$ sub-case of Prop~\ref{prop:n4mu}:
the cross-points $s_1+(t_1+2a_2)$ and $s_1+(t_1+a_2)$ give equations~\eqref{eq:cross1}
and~\eqref{eq:cross2}, and $k=2$ at $t_1-a_3$ (height~$0$) yields $\sigma_T(t_1-a_3)=-\varepsilon_1$
or a $2\varepsilon_1$ contradiction, exactly as in that sub-case---the argument goes through
unchanged since only $h_{\min}\ge 0$ and $h_{\max}=m_2$ were used, both valid here.)

If $\sigma_T(t_1-a_3)=(-1)^{m_1}\varepsilon_1$ (sub-case $|\lambda|=m_1$):
$\sigma_T(t_1+a_1+a_2)+\sigma_T(t_1+a_2)=-\varepsilon_1-(-1)^{m_1}\varepsilon_1=-(1+(-1)^{m_1})\varepsilon_1$.
If $m_1$ even: sum $=-2\varepsilon_1$; each term $\in\{-1,0,1\}$, impossible.
If $m_1$ odd: sum $=0$, so $\sigma_T(t_1+a_2)=-\sigma_T(t_1+a_1+a_2)$.
Since $t_1+a_2\in F_2$ with sign $-\varepsilon_1$ (for $m_2\ge 1$):
$\sigma_T(t_1+a_1+a_2)=\varepsilon_1$.
Then apply $k=2$ at $t_1+a_1+a_2$, using $\sigma_T(t_1+a_1)=-\varepsilon_1$ (computed above):
$\sigma_T(t_1+2a_1+a_2)+\varepsilon_1+(-\varepsilon_1)=0$: $\sigma_T(t_1+2a_1+a_2)=0$.
Continue the staircase upward; the argument is the same as in the main $m_1$ odd case.

\textit{Sub-case $m_2=0$:}
All $m_i=0$ (by closure with $\mu=0$: $0=m_4\cdot 1$, so $m_4=0$; similarly $m_3=0$ from $m_2+m_3\cdot 0=m_4\cdot 1$, giving $0=0$, consistent with any $m_3$, but $h_{\max}=0$ forces $m_3=0$).
T has both signs; same bi-infinite chain argument as in Prop~\ref{prop:n4mu}, Sub-case $m_2=0$.
\end{proof}

\medskip
\noindent\textbf{Sub-case $\mu\in\mathbb{Z}_{>0}$.}

For $\mu\in\mathbb{Z}_{>0}$ (a positive integer): the height $h(t)-\mu$ of the $a_3$-term
is an integer, and may lie in $H$.  We handle this by a modified cross-point argument.

For $t\in T$ at integer height $h\in\{0,\dots,m_2-1\}$: apply $k=1$ and $k=2$ equations.
The $a_1$-direction always preserves height, giving bi-infinite chains unless certain
height-$(h-\mu)$ elements are nonzero.  For each integer height $h'\in\{0,\dots,m_2-\mu-1\}$:
$h'$ and $h'+\mu$ are both integer heights.  The elements at height $h'$ (from $F_2$)
have known signs; the $k=3$ equation at height $h'+\mu$ gives:
$\sigma_T(t+a_1+a_2)+\sigma_T(t+a_2)+\sigma(t)+\sigma_T(t-a_3)=0$.
Here $\sigma_T(t-a_3)$ is at height $h'+\mu-\mu=h'$: a known $F_2$-sign
$(-1)^{h'}\varepsilon_1$ (if $t-a_3=t_2+h'a_2$, the face-chain element).
In general, $t-a_3$ may not equal the face-chain element; but applying the $k=1$/$k=2$
equations at height $h'$ shows any deviation from the face-chain value leads to
a bi-infinite chain.  The argument follows the same structure as Prop~\ref{prop:n4mu}:
a staircase contradiction arises from choosing a starting element whose sign disagrees
with $\varepsilon_3$ at height $m_2$, using the additional constraint that $\sigma_T(t-a_3)$
is determined (via the face-chain or a $2\varepsilon$-contradiction) at each step.

In all sub-cases of $\mu\in\mathbb{Z}_{>0}$, a magnitude-$2$ or sign contradiction is derived.

\bigskip
\noindent\textbf{Section~3: The case $n\ge 5$.}

For $n\ge 5$ with $\ge 3$ direction classes: choose the bottom edge $a_1$, set $a_2=s_3-s_2$,
define $\beta$ as above and $\mu=\beta(a_3)$.
Set $c_j=\beta(s_j-s_1)$; then $c_1=c_2=0$, $c_3=1$, $c_4=1+\mu$ (from $c_4=c_3+\beta(a_3)$).
For $j\ge 5$: $c_j\ge c_4+\beta(a_4)>c_4$ (since the polygon is strictly convex,
each subsequent vertex goes further "up" in $\beta$-direction or comes back down;
the bottom-edge choice ensures all $c_j>0$ for $j\ge 3$).

The $k=3$ equation for $t\in T$, $t\ne t_3$, at height $h(t)=h\in\{0,\dots,m_2-1\}$:
\[
  \sum_{j=1}^n \sigma_T(t+s_3-s_j)=0,
  \qquad
  h(t+s_3-s_j)=h+1-c_j.
\]
Terms $j=1,2$: height $h+1$.  Term $j=3$: height $h$.  Terms $j\ge 4$: height $h+1-c_j\le h-\mu$.
For $j=4$: height $h-\mu$.  For $j\ge 5$: height $h+1-c_j<h-\mu$ (strictly less).

By the same No-Gaps argument (Lemma~\ref{lem:nogaps}), applied now with $h_{\max}=h(t_{n})$
(the height of the last face-chain vertex; for the appropriate sub-case of $\mu$):
\begin{itemize}
\item For $\mu>0$ and $h<\mu$: height $h-\mu<0$, so all $j\ge 4$ terms vanish;
  the equation reduces to the three-term form~\eqref{eq:3term4}.
\item For $\mu>0$ and $h\ge\mu$: the $j=4$ term at height $h-\mu$ is non-integer
  (for $\mu$ non-integer) or is handled by the extended No-Gaps argument;
  similarly for $j\ge 5$.
\item For $\mu\le 0$: heights $h+1-c_j$ for $j\ge 4$ satisfy $h+1-c_j\le h-\mu\le h$
  (since $-\mu\ge 0$ and $c_j\ge 1+\mu>0$); the terms vanish by Lemma~\ref{lem:betamax4neg}
  or Lemma~\ref{lem:nogaps} applied to the heights $h-|\mu|,\dots$
\end{itemize}

In all cases, for $h\in\{0,\dots,m_2-1\}$ the $k=3$ equation reduces to:
\[
  \sigma_T(t+a_1+a_2)+\sigma_T(t+a_2)+\sigma(t)=0
\]
(the same three-term form as for $n=3$), and the $j\ge 4$ terms all vanish.

The heights $h(t_1)=h(t_2)=0$ and $h(t_3)=m_2$ still hold.
The closure equation for $n\ge 5$ gives a more complex relation on the $m_i$, but
Lemma~\ref{lem:betamin4} (with $c_j>0$ for $j\ge 3$) and the height-maximum lemma
(using the $k=n$ equation analogously to Lemma~\ref{lem:betamax4mu}/\ref{lem:betamax4neg})
still guarantee $h_{\min}=0$ and a well-defined $h_{\max}$.

Moreover, $h_{\max}\ge m_2$ (since $t_3\in T$ at height $m_2$), and the uniqueness of $t_3$
at height $m_2$ (among elements relevant to the three-term staircase)
follows from the three-term equation at height $m_2-1$ combined with $h_{\max}$ considerations.

With the three-term staircase in place, the starting-element argument of Prop~\ref{prop:n4mu}
(applied to the sub-system $\{s_1,s_2,s_3\}$ of the face chains $F_1$ and $F_2$)
yields the sign contradiction at height $m_2$ for the appropriate sub-case of $m_1$.
(Note: $m_1=m_2$ need not hold for $n\ge 5$; the starting-element choice depends only on the
parity of $m_1$, which suffices.)

\bigskip
\noindent\textbf{Conclusion.}
For any $n\ge 3$ and any convex polygon $P$ with $\ge 3$ edge-direction classes,
the extremal equality $|\operatorname{supp}(\sigma)|=n$ with both signs in~$T$ leads to a contradiction:
\begin{itemize}
  \item $n=3$: Proposition~\ref{prop:n3}.
  \item $n=4$, $\mu\notin\mathbb{Z}_{\ge 0}$, $\mu\ne 0$: Proposition~\ref{prop:n4mu}.
  \item $n=4$, $\mu=0$: Proposition~\ref{prop:n4mu0}.
  \item $n=4$, $\mu\in\mathbb{Z}_{>0}$: the modified cross-point argument described above.
  \item $n\ge 5$: Section~3 reduction to the three-term staircase.
\end{itemize}
$\blacksquare$

\end{document}
